\section{Introduction}

3D Scene Graphs turn reconstructed scenes into structured object-relation
representations that can support spatial reasoning, embodied AI, scene
alignment, visual grounding, and downstream decision making
\cite{wald2020learning3dssg,sarkar2023sgaligner,xie2024sgpgm,liu2026viewongraph}.
For these uses, relation prediction must be reliable in the physical scene, not
only plausible at the category or language level. A graph edge such as
\texttt{standing on}, \texttt{next to}, or \texttt{higher than} is useful only
if it describes the actual subject-object pair in 3D.

We study a failure mode in which semantic relation predictors rank plausible
relation labels that are geometrically inconsistent for the actual object pair.
A relation can sound plausible for two object categories while being
contradicted by their contact evidence, distance, or vertical arrangement in
the reconstructed scene. This mismatch matters more as 3D Scene Graphs become
interfaces for language-facing, open-vocabulary, and embodied systems
\cite{koch2024open3dsg,hou2025fross,zhang2025openfungraph,wang2025octreegraph,saxena2025zing3d,yu2025openworld3dsg,madhavaram2026vizor}.

The failure is not simply that existing 3DSSG methods ignore geometry. Prior
work already uses edge features, point-cloud structure, spatial knowledge,
multimodal semantics, and geometric fusion
\cite{zhang2021sggpoint,feng2023smka,wang2023vlsat,koch2024sgrec3d,chen2024ccl3dsgg}.
The gap is more specific: semantic relation confidence is not necessarily a
calibrated estimate of relation-level physical consistency. A high relation
score can behave as a semantic plausibility score without indicating whether
the same subject-object pair satisfies the relation in 3D.

This motivates a \method{} framework. The framework treats an existing
predictor as a relation source, standardizes its outputs into
identity-preserving prediction rows, joins object-pair geometry evidence by
subject and object instance identifiers, estimates \pgeom{}, and evaluates
relation rankings with both exact-label recall and geometric violation. This
design is necessary because hard physical rules and probabilistic reliability
scores answer different questions: rule-verified variants diagnose strict
consistency, while calibrated variants expose a usable recall-violation
operating point. Figure~\ref{fig:framework} summarizes this path.

We evaluate the framework on geometry-checkable relation families:
\texttt{support\_contact}, \texttt{proximity}, and
\texttt{relative\_vertical}. VL-SAT is the primary reproduced closed-set
relation source, and Open3DSG is used as measured second-source evidence within
the same H001-family scope \cite{wang2023vlsat,koch2024open3dsg}. Controls test
whether the effect can be explained by geometry-only ranking, a generic
distance heuristic, shuffled geometry, or wrong-pair geometry. GT-positive and
counterfactual checks test the geometry-validity signal independently from
prediction re-ranking, while qualitative failure analysis shows how semantic
plausibility diverges from physical consistency.

Our contributions are threefold:
\begin{enumerate}
    \item We identify and formalize a relation-level reliability failure in 3D
    Scene Graph prediction: semantic relation confidence can rank plausible
    predicates that are physically inconsistent because it is not calibrated to
    object-pair geometry.
    \item We introduce a \method{} framework that standardizes prediction rows,
    joins identity-preserving 3D evidence, estimates \pgeom{}, and exposes
    probabilistic, rule-verified, and family-specific operating points.
    \item We define a recall-violation evaluation protocol for
    geometry-checkable relation families, including exact-label \rAt{K},
    \vAt{K}, GT-positive/counterfactual verifier evaluation, and geometry
    identity controls.
\end{enumerate}

We validate these contributions with Docker-reproducible evidence on VL-SAT and
Open3DSG, including nontriviality controls, second-source metrics, and failure
analysis. The claim is intentionally scoped: H001 measures and improves
relation reliability for geometry-checkable 3DSSG families. It does not claim
general open-vocabulary 3D Scene Graph generation improvement,
arbitrary-source generality, or guaranteed physical correctness.

